\subsection*{From Expressions to Value Numbers}

The current value-numbering implementation follows the syntax tree. It traverses the abstract syntax tree and assigns a value number whenever a node contains an expression.

The method of generating value numbers is simple. For each expression, it takes the expression as the input, and returns corresponding value number.

More specifically, it:

\begin{itemize}[leftmargin=1.6em]
\item Generates a signature from an expression.
\item If the signature already exists, return the associated value number by looking up the value number table.
\item Or else, it associates the signature with a new value number in the value number table, and return it.
\end{itemize}

For simplicity, the current implementation represents an expression signature as a string. A hash-based representation could replace it later, as in many value-numbering implementations.

In addition, to handle the symbol definition and symbol reference, it follows the following rules:

\begin{itemize}[leftmargin=1.6em]
\item When an expression is to initialize a symbol (variable or partial type), the symbol is assigned the same value number of the expression. More specifically, it associates the signature of the \emph{scoped} symbol name with the value number in the value number table.
\item When an expression references a symbol, it looks up its corresponding signature in the value number table, and return it.
\end{itemize}

Note the symbol must use its \emph{scoped name} for signature generation, since identical names are allowed to be defined in different scopes.
